\documentclass[10pt,a4paper]{amsart}
\usepackage[margin=19mm]{geometry}
\usepackage{amsmath,amssymb,mathtools}
\usepackage[scr=rsfs,cal=euler]{mathalfa}
\usepackage{xfrac}
\usepackage[unicode,hidelinks,bookmarks=false]{hyperref}
\newcommand{\ie}{i.e.\ }
\newcommand{\cf}{cf.\ }
\newcommand{\resp}{respectively\ }
\newcommand{\Subsection}{\subsection}
\providecommand{\nobreakdash}{\nobreak\hbox{-}}
\setlength{\emergencystretch}{2em}
\multlinegap0pt
\usepackage{color}
\usepackage{xypic}
\usepackage{thmtools}
\numberwithin{equation}{section}
\newtheorem{theorem}{Theorem}[section]
\newtheorem{prop}[theorem]{Proposition}
\newtheorem{lemma}[theorem]{Lemma}
\newtheorem{cor}[theorem]{Corollary}
\newtheorem{claim}{Claim}
\newtheorem{conjecture}{Conjecture}
\newtheorem{definition}[theorem]{Definition}
\newtheorem{exa}[theorem]{Example}
\newtheorem{rem}[theorem]{Remark}
\newtheorem{que}[theorem]{Question}
\newcommand{\st}{\ :\ }
\newcommand{\A}{\mathbf{A}}
\newcommand{\B}{\mathbf{B}}
\newcommand{\C}{\mathcal{C}}
\newcommand{\D}{\mathbf{D}}
\newcommand{\F}{\mathbf{F}}
\newcommand{\K}{\mathbf{K}}
\renewcommand{\L}{\mathbf{L}}
\newcommand{\M}{\mathbf{M}}
\newcommand{\N}{\mathbb{N}}
\newcommand{\R}{\mathbf{R}}
\renewcommand{\S}{\mathbf{S}}
\newcommand{\U}{\mathbf{U}}
\newcommand{\V}{\mathcal{V}}
\newcommand{\Z}{\mathbb{Z}}
\newcommand{\Aut}{\mathrm{Aut}}
\newcommand{\Ann}{\mathrm{Ann}}
\newcommand{\Clo}{\mathrm{Clo}}
\newcommand{\Con}{\mathrm{Con}}
\newcommand{\GL}{\mathrm{GL}}
\newcommand{\Pol}{\mathrm{Pol}}
\newcommand{\rk}{\mathrm{rk}}
\newcommand{\Hom}{\mathrm{Hom}}
\newcommand{\Cg}{\mathrm{Cg}}
\newcommand{\SMP}{\text{SMP}}
\newcommand{\blue}[1]{{\color{blue} #1}}
\newcommand{\pmbegin}{\color{blue}}
\newcommand{\pmend}{\color{black}}
\newcommand{\red}{\color{red}}
\newcommand{\redend}{\color{black}}
\begin{document}
\title[Structure and Complexity of 2-Nilpotent Mal'cev Algebras]{Structure and Complexity of 2-Nilpotent Mal'cev Algebras}
\author{Patrick Wynne}
\date{}
\maketitle
\noindent\emph{Source and licence.} Structure and Complexity of 2-Nilpotent Mal'cev Algebras. Version arXiv:2608.18917v1 (CC BY 4.0). This material is distributed under the Creative Commons Attribution 4.0 International licence, http://creativecommons.org/licenses/by/4.0/, as recorded in the arXiv OAI licence metadata for this exact version and shown on the abstract page https://arxiv.org/abs/2608.18917v1. Full rights notice: licenses/arxiv-2608.18917-CC-BY-4.0.txt.\par\smallskip
\noindent\textit{Editorial scope.} Counted unit: the original \texttt{theorem} environment \texttt{thm:clonegens} at source lines 686--696 together with its complete proof at lines 698--723; the delivered document keeps source lines 391--724 so that every referenced label and every citation resolves inside it. Native editable source is the paper's own concatenated TeX; the AMS wrapper is editorial formatting only. No publisher class, upstream script or unrelated artwork is redistributed.
\section{Central extensions} 
\label{sec:centralextensions}

Let $\V$ be a congruence modular variety over the signature $F$. 
Then $\V$ has a ternary \emph{difference term} $d$ 
  such that for all $\A\in \V$ and for all $x,y\in A$,
  \[ d^\A(x,x,y) = y  \quad \text{and} \quad  d^\A(y,x,x) \equiv y \mod [\alpha,\alpha], \]
  where
  \[ \alpha = \Cg_\A(x,y) := \bigcap \{\theta \in \Con(\A) \st (x,y) \in \theta\},\]
  is the smallest congruence relating $x$ and $y$ and $[\alpha,\alpha]$ is the commutator 
  of $\alpha$ with itself. 

Note that if $\L \in \V$ is abelian then $[\theta, \theta] = 0$ for all $\theta \in \Con(\A)$.
So $d^\L(x,x,y) = y$ and $d^\L(x,y,y) = x$ for all $x,y \in L.$ 
That is, $d$ interprets as a Mal'cev operation on $\L$. 
 
Let $\V$ be a congruence modular variety with signature $F$ and difference term $d$. 
Let $\L, \U \in \V$. Suppose that $\L$ is an abelian algebra with associated group 
  $\hat{\L} = (L, +, -, 0)$ 
  so that $d^\L(x,y,z) = x-y+z$ for all $x,y,z \in L$. 
See \cite{FM:CTC}~Chapter~5 for a detailed explanation of the abelian group associated 
  to an abelian Mal'cev algebra.
Suppose for each operation symbol $f \in F$ we are given a map $\hat{f} \colon U^k \rightarrow L$ 
  where $k$ is the arity of $f$. 
Let $\hat{F} := \{\hat{f} \st f \in F\}$. We define a new algebra $\L \otimes^{\hat{F}} \U$ 
  with universe $L \times U$ and basic operations 
  \[f^{\L \otimes^{\hat{F}} \U} \colon (L \times U)^k \rightarrow L \times U, \, (\ell, u) \mapsto (f^\L(\ell) + \hat{f}(u), f^\U(u)) \quad \text{ for } f \in F. \]
Here and throughout we slightly abuse notation and let $\ell = (\ell_1,\dots,\ell_k)$ 
  and $u = (u_1,\dots,u_k)$ for $( (\ell_1,u_1),\dots,(\ell_k,u_k) ) \in (L\times U)^k$.
We note that $\L \otimes^{\hat{F}} \U$ may not be in the variety $\V$. 
 
In \cite{FM:CTC} Freese and McKenzie show that in a congruence modular variety $\V$ every algebra $\A$
  that is isomorphic to $\U$ modulo a central congruence $\alpha$ is of the form $\L\otimes\U$ 
  for some abelian $\L\in \V$.
 
\begin{theorem}[\cite{FM:CTC} Theorems 7.1 \& 7.2] 
  Let $\V$ be a congruence modular variety with signature $F$ and let $\A \in \V$. 
  For $\zeta_\A$ the center of $\A$, let $\U = \A/\zeta_\A$.
  \begin{enumerate}
    \item  
      There is an abelian algebra $\L \in \V$ and a system $\hat{F}$ such that $\A \cong \L \otimes^{\hat{F}}\U$. 
      Moreover, the center of $\L \otimes^{\hat{F}} \U$ is the kernel of the projection onto $U$. 
    \item 
      $\A$ is $2$-nilpotent if and only if $\A \cong \L \otimes^{\hat{F}}\U$ for some system $\hat{F}$ 
        and for $\L, \U$ abelian algebras in $\V$. 
  \end{enumerate}
\end{theorem}
 
We call $\L \otimes^{\hat{F}} \U$ a \emph{central extension} of $\L$ by $\U$. 
We note that $\L \otimes^{\hat{F}} \U$ is sometimes called a wreath product or a 
  semidirect product in the literature. 
We will usually suppress the superscript $\hat{F}$ and simply write $\L \otimes \U$. 
Since $\L$ is an abelian Mal'cev algebra, it is straightforward to show by induction on complexity of terms 
  that for each $k \in \N$, every $k$-ary term $t$ over $F$ induces a term function $t^{\L \otimes \U}$ of $\L \otimes \U$ of the form
  \[t^{\L \otimes \U} \colon (L \times U)^k \rightarrow L \times U, \, (\ell, u) \mapsto (t^\L(\ell) + \hat{t}(u), t^\U(u))\]
  for some $\hat{t}\colon U^k \rightarrow L$. 
We will often abuse notation and write $t^{\L \otimes \U} = (t^\L, t^\U) + \hat{t}$. 
We call $\hat{t} \colon U^k \rightarrow L$ the distortion of the term function $t^{\L \otimes \U}$. 
More generally for any $k$-ary term function $t$ on $\L \otimes \U$ and for any 
  $e \colon U^k \rightarrow L$, we abuse notation and denote by $t + e$ the $k$-ary 
  function on $L \times U$ defined by
  \[(\ell, u) \mapsto (t^\L(\ell)+\hat{t}(u) + e(u), t^\U(u)).\]
 
\begin{definition} \label{def:elprime} 
  For $\V$ a congruence modular variety and $\L \in \V$ an abelian algebra we define a new algebra $\L'$ 
    of signature $F$ with universe $L$ and basic operations 
    \[f^{\L'}(\ell_1, \ldots, \ell_k) := f^\L(\ell_1, \ldots, \ell_k) - f^\L(0, \ldots, 0)\]
    for $f \in F$ of arity $k$ and for $0 \in L$.
\end{definition}

Note that $\L' = \L$ if and only if $\{0\}$ is a subuniverse of $\L$. 
Moreover, for $t$ an $F$-term of arity $k$,  by induction on the complexity of terms, we see that 
  \[t^{\L'}(\ell_1, \ldots,\ell_k) = t^\L(\ell_1, \ldots, \ell_k) - t^\L(0, \ldots, 0).\] 
We note that $\L$ and $\L'$ are polynomially equivalent and have the same idempotent term functions. 
Moreover, we see that $\L' \in \V$ as follows.
Let $s$ and $t$ be $F$-terms and suppose that $s^\L = t^\L$ is a $k$-ary identity of $\L$.
Then $s^\L(0, \dots, 0) = t^\L(0, \dots, 0)$ for $0 \in L$,
  and so $s^\L(x) - s^\L(0, \dots, 0) = t^\L(x) - t^\L(0, \dots,0)$ for all $x \in L^k$. 
That is, $s^{\L'} = t^{\L'}$ is an identity of $\L'$.
So $\L'$ satisfies all identities satisfied in $\L$, and hence $\L' \in \V(\L)$. 
In particular, since $\L$ is abelian, $\L'$ is also abelian with Mal'cev term 
  $d^{\L'}(x,y,z) = d^\L(x,y,z) = x-y+z$.
We will often consider the algebra $\L'$ expanded by $0$, which we denote by $\L'_0$.
 
For $f\in F$ define $f'(u) := \hat{f}(u)+f^\L(0,\dots,0)$. 
Let $F' := \{f'\st f\in F\}$. 
Then
  \[ \L\otimes^{\hat{F}}\U = \L'\otimes^{F'}\U. \] 
 
The easiest example of a central extension of $\L$ by $\U$ is the direct product, $\L \times \U$. 
In this case, $\hat{f} = 0$ for every $f \in F$. 
 
For clones $C$ and $D$, we say that a map $\phi \colon C \rightarrow D$ is a \emph{clone homomorphism} if $\phi$
  \begin{enumerate}
    \item preserves arities, i.e. the arity of $f$ equals the arity of $\phi(f)$ for all $f \in C$,
    \item preserves projections, i.e. maps projections in $C$ to the corresponding projections in $D$,
    \item and preserves generalized compositions, i.e. \[\phi(f(g_1,\ldots,g_k)) = \phi(f)(\phi(g_1), \ldots, \phi(g_k)) \] 
      for all $f, g_1, \ldots, g_k \in C$ of the appropriate arities. 
  \end{enumerate}
We note that if $\phi \colon C \rightarrow D$ is a bijective clone homomorphism then $\phi^{-1} \colon D \rightarrow C$ 
  is a clone homomorphism.
In this case we call $\phi$ a clone isomorphism.
 
The following lemmas relate arbitrary central extensions to the direct product via a clone homomorphism. 
 
\begin{lemma} \label{lem:clonehom}
  Let $\V$ be a congruence modular variety and let $\L \otimes \U \in \V$ be a central extension 
    of an abelian $\L$ by $\U$. Let $\L'$ be as in Definition~\ref{def:elprime}. 
  The map
    \[ \xi\colon\Clo(\L\otimes\U) \to \Clo(\L'\times\U),\ f^{\L\otimes \U} \mapsto f^{\L'\times\U}, \]
    is a surjective clone homomorphism. 
\end{lemma}

\begin{proof}  
  Note that $\xi$ is well-defined since for $k$-ary $F$-terms $f,g$ with 
    $f^{\L\otimes \U} = g^{\L\otimes \U}$, we have $f^\U=g^\U$ and
    \[ f^\L(\ell)+\hat{f}(u) = g^\L(\ell)+\hat{g}(u) \text{ for all } \ell\in L^k, u\in U^k. \]
  Setting $\ell = (0,\dots,0)$ in the latter, we obtain
    \[ f^\L(0,\dots,0)+\hat{f}(u) = g^\L(0,\dots,0) +\hat{g}(u) \text{ for all } u\in U^k. \]
  Taking the difference in $\L$ of the two equations above we obtain
    \[ f^\L(\ell) - f^\L(0,\dots,0) = g^\L(\ell)-g^\L(0,\dots,0) \text{ for all } \ell\in L^k. \]
  Thus $f$ and $g$ induce the same function on $\L'$ and $\xi$ is well-defined.
 
  From its definition it is immediate that $\xi$ preserves arities, projections, and the composition of functions.  
  Moreover, let $f$ be a term in the signature of $\V$ that induces the term function 
    $f^{\L' \times \U}$ on $\L' \times \U$. 
  Then $f$ induces a term function $f^{\L \otimes \U}$ on $\L \otimes \U$, and 
    $\xi(f^{\L \otimes \U}) = f^{\L' \times \U}$.
  Hence $\xi$ is surjective.
\end{proof}

Note that  for any central extension $\L \otimes \U$, the projection 
  $\pi_U \colon \L \otimes \U \rightarrow \U$ is a homomorphism.

\begin{lemma} \label{lem:injective} 
  Let $\L \otimes \U$  be as in Lemma~\ref{lem:clonehom} and let $\L'$ be as in Definition~\ref{def:elprime}.
  The projection $\pi_L \colon \L \otimes \U \rightarrow \L'$ is a homomorphism if and only if 
    $\L \otimes \U = \L' \times \U$. 
\end{lemma}

\begin{proof}
  Suppose that $\pi_L \colon \L \otimes \U \rightarrow \L'$ is a homomorphism.
  Let $s$  be an $F$-term and let $\ell \in L^k$ and $u \in U^k$.
  Then 
    \[ \pi_L(s^{\L \otimes \U}(\ell,u))  = \pi_L(s^\L(\ell) + \hat{s}(u), s^\U(u))  = s^\L(\ell) + \hat{s}(u). \]
  Since $\pi_L$ is a homomorphism, 
    \[\pi_L(s^{\L \otimes \U}(\ell,u)) = s^{\L'}(\pi_L(\ell, u)) = s^{\L'}(\ell)= s^\L(\ell) - s^\L(0, \ldots,0).\]
  So $\hat{s}(u) = -s^\L(0, \ldots,0)$.
  Therefore,
    \[ s^{\L \otimes \U} = (s^\L + \hat{s}, s^\U) = (s^\L - s^\L(0, \ldots, 0), s^\U) = s^{\L' \times \U}. \]
  Hence $\L \otimes \U = \L' \times \U$.
  Conversely if $\L \otimes \U = \L' \times \U$ then clearly $\pi_L$ is a homomorphism. 
\end{proof}

The following theorem will show that $\L \otimes \U$ is abelian when $\xi$ is injective and $\U$ is abelian. 

\begin{theorem}\cite[Theorem 9.3]{ALV:III} \label{ALV:clone} 
  Suppose $\A$ and $\B$ are two algebras of signature $F$. 
  \begin{enumerate}
    \item 
      $\B \in \V(\A)$ if and only if there is a clone homomorphism 
        \[\phi \colon \Clo(\A) \rightarrow \Clo(\B)\]
        with $\phi(f^\A) = f^\B$ for all $f \in F$.
    \item 
      $\V(\A) = \V(\B)$ if and only if there exists a clone isomorphism 
        \[\phi \colon \Clo(\A) \rightarrow \Clo(\B)\]
        with  $\phi(f^\A) = f^\B$ for all $f \in F$.
  \end{enumerate}
\end{theorem}

\begin{cor} 
  Let $\xi$ and $\L \otimes \U$  be as in Lemma~\ref{lem:clonehom} and suppose that $\xi$ is injective. 
  For $k \in \N$, if $\U$ is $k$-(super)nilpotent, then $\L \otimes \U$ is $k$-(super)nilpotent.
  In particular if $\U$ is abelian, then $\L \otimes \U$ is abelian. 
\end{cor}
\begin{proof}
  Let $\L'$ be as in Lemma~\ref{lem:clonehom}.
  Since $\xi$ is a injective and surjective, 
    $\xi \colon \Clo(\L \otimes \U) \rightarrow \Clo(\L' \times \U)$ is a clone isomorphism.
  By Theorem~\ref{ALV:clone}, $\V(\L \otimes \U) = \V(\L' \times \U)$.
  In particular, $\L \otimes \U \in \mathbb{HSP}(\L' \times \U)$.
  Since quotients, subalgebras, and products of $k$-(super)nilpotent algebras in a congruence modular variety 
  remain $k$-(super)nipotent, the result follows. 
\end{proof}

\begin{exa} 
  We show that for abelian $\L$ and $\U$, the central extension $\L \otimes \U$ may be abelian even if $\xi$ is not injective. % 
  Let $\L = (\Z_2, +) = \U$. 
  Define $\L \otimes \U = (L \times U, p)$ 
    where 
    \[p^{\L \otimes \U}((\ell_1, u_1) , (\ell_2, u_2))= (\ell_1 + \ell_2 + \hat{p}(u_1,u_2), u_1 + u_2)\]
    with 
    \[ \hat{p} \colon U^2 \rightarrow L, \quad  (u_1, u_2) \mapsto 
      \begin{cases} 
        1  & \text{ if } u_1 = u_2 = 1, \\ 
        0 & \text{ else.} 
      \end{cases}\]
  Then $\L \otimes \U$ is a cyclic abelian group, and hence isomorphic to $(\Z_4, +)$.

  Note that $\xi \colon \Clo(\L \otimes \U) \rightarrow \Clo(\L \times \U)$ is not injective here. 
  So for abelian $\L$ and $\U$, we see that $\L \otimes \U$ may be abelian even if  $\L \otimes \U \ncong \L' \times \U$.
\end{exa}

To capture the situation when $\xi$ is not injective we define the \emph{difference clonoid} of $\L\otimes\U$ as follows.
 
\begin{definition} \label{def:differenceclonoid} 
  Let $\V$ be a congruence modular variety and let $\L \otimes \U$ be a central extension of the abelian $\L$ by $\U$.
  Let  $\L'$ be as in Definition~\ref{def:elprime}.
  For $0 \in L$ we define the \emph{difference clonoid} from $\U$ to $\L'_0$, the algebra $\L'$ expanded by $0$, as
    \[ D(\L\otimes\U) := \{ e\colon U^k\to L \st k\in\N, \,\, x_1+e \in\Clo_k(\L\otimes\U) \}. \]
\end{definition}

We prove that $D(\L\otimes\U)$ is indeed a clonoid in the following.
 
\begin{lemma} \label{lem:diff} 
  Let $\V$ be a congruence modular variety and let $\L \otimes \U \in \V$ be a central extension of $\L$ by $\U$. 
  Let $0 \in L$ and let $\xi$ be as in Lemma~\ref{lem:clonehom}.  
  Let  $\L'$ be as in Definition~\ref{def:elprime}.
  \begin{enumerate}
    \item  \label{it:addon}
      For $k\in\N$ and $e\colon U^k\to L$ the following are equivalent:
      \begin{enumerate} 
        \item $e\in D(\L\otimes\U)$;
        \item $f+e \in\Clo_k(\L\otimes\U)$ for all $f\in\Clo_k(\L\otimes\U)$;
        \item $f+e \in\Clo_k(\L\otimes\U)$ for some $f\in\Clo_k(\L\otimes\U)$;
        \item $(f, f+e) \in \ker(\xi)$ for all $f \in \Clo_k(\L \otimes \U)$.
      \end{enumerate}   
    \item \label{it:clonoid}
      $D(\L\otimes\U)$ is a clonoid from $\U$ to $\L'$ expanded with $0$, in particular to $(L,+,-,0)$.
    \item \label{it:zero} 
      $\xi$ is a clone isomorphism if and only if
        \[D(\L\otimes\U) = \{ f \colon U^k \rightarrow L \st k \in \N, \,  f(x) = 0 \text{ for all } x \}.\]
  \end{enumerate}
\end{lemma}

\begin{proof}
  Let $D := D(\L\otimes\U)$, and let $d$ be a difference term for $\V$.
  Since $\L$ is abelian, $d^\L(x,y,z) = x-y+z$.
  Then \[d^{\L \otimes \U}((\ell_1,u_1),(\ell_2,u_2),(\ell_3,u_3)) = (\ell_1-\ell_2+\ell_3 + \hat{d}(u_1,u_2,u_3), d^\U(u_1,u_2,u_3)).\]
  Since $d$ is a difference term, $\hat{d}(u_1,u_1,u_2) = 0 = \hat{d}(u_2,u_1,u_1)$ for all $u_1,u_2\in U$.

  \eqref{it:addon}
    Let $f,g\in\Clo_k(\L\otimes\U)$ such that $f+e\in\Clo_k(\L\otimes\U)$. 
    We claim 
      \begin{equation} \label{eq:dfe}
        d^{\L\otimes\U}(f+e,f,g) = g+e.
      \end{equation}

  To show~\eqref{eq:dfe} let $\ell\in L^k,u\in U^k$ and consider 
    \begin{align*}
      d^{\L\otimes\U}&(f+e,f,g) (\ell,u)  \\
      = & (d^\L(f^\L(\ell)+\hat{f}(u)+e(u),f^\L(\ell)+\hat{f}(u),g^\L(\ell)+\hat{g}(u)) \\
      & \quad \quad +\underbrace{\hat{d}(f^\U(u),f^\U(u),g^\U(u))}_{=0},  d^\U(f^\U(u),f^\U(u),g^\U(u)) ) \\
      = & (g^\L(\ell)+\hat{g}(u)+e(u), g^\U(u)) \\
      =  & (g+e)(\ell,u). 
    \end{align*}
  All equivalences in \eqref{it:addon} now follow from~\eqref{eq:dfe}.


  \eqref{it:clonoid}
    From~\eqref{it:addon} it follows that $D$ is closed under $+,-,0$ on $L$. 
    More generally, let $f\in F$ be $k$-ary  and
      let $e_1,\dots,e_k\in D$ be $n$-ary. 
    Let $x_i = (\ell_i, u_i) \in L \times U$ for $1 \le i \le n$ and let $u = (u_1, \dots, u_n) \in U^n$.
    Since $\L$ is abelian, we have
    \begin{align*}
      f^{\L\otimes \U} & (x_1+e_1(u), \dots, x_1 +e_k(u)) \\
      &= ( f^\L(\ell_1 + e_1(u), \dots \ell_1 + e_k(u)) + \hat{f}(u_1, \dots, u_1) , f^\U(u_1, \dots, u_1) ) \\
      & =  ( f^\L(\ell_1 - 0 + e_1(u), \dots , \ell_1 - 0 + e_k(u)) + \hat{f}(u_1, \dots, u_1) , \\
      & \quad \quad f^\U(u_1, \dots, u_1) ) \\ 
      & = ( f^\L(\ell_1, \dots, \ell_1) - f^\L(0, \dots, 0)  \\
      & \hspace{1in}+ f^\L(e_1(u), \dots, e_k(u)) + \hat{f}(u_1, \dots, u_1), \\
      & \quad \quad \quad \quad \quad \quad \quad f^\U(u_1, \dots, u_1))\\
      & =  f^{\L\otimes \U} (x_1, \dots, x_1) - f^\L(0, \dots, 0) + f^\L(e_1(u), \dots, e_k(u)) \\
      &  = f^{\L\otimes \U} (x_1, \dots, x_1)  +f^{\L'} (e_1(u), \dots ,e_k(u)).
    \end{align*}
  Hence by \eqref{it:addon},
    $f^{\L'} (e_1,...,e_k) \in D$.
  Thus $\Clo(\L')D \subseteq D$.

  To see $D\Clo(\U) \subseteq D$, let $e\in D$ be $k$-ary, and let $f_1,\dots,f_k\in\Clo(\L\otimes\U)$ be $n$-ary.
  Then $x_1+e$ and consequently $(x_1+e)(f_1,\dots,f_k)$ is a term function of $\L\otimes\U$.
  For $\ell\in L^n,u\in U^n$ we see that
    \begin{align*}
      (x_1+e)(f_1,\dots,f_k)(\ell,u) & = (f^\L_1(\ell)+ e(f_1^\U(u),\dots,f^\U_k(u)), f_1^\U(u) ) \\
      & = (f_1+ e(f_1^\U,\dots,f^\U_k)) (\ell,u).           
    \end{align*}
  From~\eqref{it:addon} it follows that $e(f_1^\U,\dots,f^\U_k)  \in D$ and $D$ is a clonoid.
 
  \eqref{it:zero} 
    Follows from \eqref{it:addon}.
\end{proof}

Next we determine a set of generators for the term clone of $\L \otimes \U$. 


\begin{theorem} \label{thm:clonegens}
  Let $\L,\U$ be algebras in a congruence modular variety $\V$ with difference term $d$, 
    let $\L\otimes\U$ be a central extension of $\L$ by $\U$ and let $\xi$ be the clone 
    homomorphism from $\Clo(\L\otimes\U)$ onto $\Clo(\L'\times\U)$ from Lemma~\ref{lem:clonehom}.
 
  Let $G \subseteq \Clo(\L\otimes\U)$ such that $\xi(G)$ generates $\Clo(\L'\times\U)$.
  Let $0 \in L$ and let $E$ be a generating set of
    the clonoid $D(\L\otimes\U)$ from $\U$ to $(L, +, -, 0)$.
  Then
    \[ \Clo(\L\otimes\U) = \langle \{d^{\L\otimes\U}\} \cup G \cup (x_1+E) \rangle. \]
\end{theorem}

\begin{proof}
  The inclusion $\supseteq$ is clear since $S :=  \{d^{\L\otimes\U}\} \cup G \cup (x_1+E)$ 
    is a set of term functions of $\L\otimes\U$.
  For the converse inclusion let $t\in\Clo(\L\otimes\U)$. 
  By assumption we have $s\in\langle G\rangle$ such that $\xi(s) = \xi(t)$. 
  Hence by Lemma~\ref{lem:diff} we have $e\in D(\L\otimes\U)$ such that
    \[ t=s+e = d^{\L\otimes\U}(x_1+e,x_1,s). \]
  Thus $t\in\langle S\rangle$ follows once we have proved that
    \begin{equation} \label{eq:x1+D} 
      x_1+D(\L\otimes\U) \subseteq \langle S\rangle.
    \end{equation}
  For this we show that
    \[ D := \{e\in D(\L\otimes\U) \st x_1+e \in \langle S\rangle \} \]
    is a clonoid from $\U$ to $(L,+,-,0)$. 
  Clearly $0\in D$. For $e_1,e_2\in D$ we see from
    \[ d^{\L\otimes\U}(x_1+e_1, x_1, x_1+e_2) = x_1+(e_1+e_2),\quad d^{\L\otimes\U}(x_1,x_1+e_1,x_1) = x_1-e_1 \]
    that $D$ is closed under $+,-$ on the outside.

  Next let $e\in D$ be $k$-ary and let $p_1^\U,\dots,p_k^\U$ be $n$-ary term functions of $\U$. 
  Then there exist $n$-ary functions $q_1,\dots,q_k\in\langle G\rangle$ such that
    \[q_{i}(\ell,u) = (p_{i}^\L(\ell)+ \hat{p}_{i}(u), p_{i}^\U(u)) \text{ for all } i\in [k],\ell\in L^n,u\in U^n.\]
  Moreover, as in the proof of Lemma~\ref{lem:diff} item \ref{it:clonoid}, we see that
    \[ (x_1+e)(q_1,\dots,q_k) = q_1+e(q_{1},\dots,q_{k}) = q_1+e(p_1^\U,\dots,p_k^\U). \]
  Hence $D \Clo(\U) \subseteq D$. Thus $D$ is a clonoid from  $\U$ to $(L,+,-,0)$. Since $E\subseteq D$ by definition,
    $D=D(\L\otimes\U)$ and~\eqref{eq:x1+D} is proved. The result follows.
\end{proof}


\begin{thebibliography}{99}
\bibitem[ALV:III]{ALV:III} Ralph~S. Freese, Ralph~N. McKenzie, George~F. McNulty, and Walter~F. Taylor. \newblock {\em Algebras, lattices, varieties. {V}ol. {III}}, volume 269 of {\em Mathematical Surveys and Monographs}. \newblock American Mathematical Society, Providence, RI, 2022.
\bibitem[FM:CTC]{FM:CTC} R.~Freese and R.~N. McKenzie. \newblock {\em Commutator Theory for Congruence Modular Varieties}, volume 125 of {\em London Math. Soc. Lecture Note Ser.} \newblock Cambridge University Press, 1987. \newblock Available from \verb+http://math.hawaii.edu/~ralph/Commutator/comm.pdf+.
\end{thebibliography}
\end{document}
